\section{Two twists and a finite bimodule fiber}
\label{branch:section}

We now construct the finite bimodule and the stable map required by
Proposition~\ref{conv:proposition}.  The two twists will act differently on every
positive self-extension of $X$, while agreeing on its identity.  The fiber
of their maps to $\mathcal C[3]$ turns this distinction into the comparison
isomorphisms used in that proposition.

\subsection{A finite representative and an exact sequence}

Write $R=A^e$.  Lemma~\ref{cone:finite} provides the finite bimodule
\[
 C_1=\coker(K_2\longrightarrow K_1),
 \qquad \mathcal C=C_1[1]
\]
in the stable category of $R$.  It is projective on both $A$-sides.  To fix
actual finite representatives for its shifts, for a finite left $R$-module
$M$ use the injection
\begin{equation}
 \label{branch:finite-embedding}
 j_M:M\longrightarrow I_R(M):=\Hom_{\kk}(R,M),
 \qquad j_M(m)(t)=tm,
 \qquad (r\varphi)(t)=\varphi(tr).
\end{equation}
Evaluation at $1$ proves injectivity.  Since $R$ is symmetric,
$I_R(M)\cong DR\otimes_{\kk}M\cong R\otimes_{\kk}M$ as left $R$-modules,
so this is a finite projective module.  Put
$\Sigma_RM=\coker j_M$.  When $M$ is projective on each $A$-side, the
sequence defining $\Sigma_RM$ splits on each side: $A$ is self-injective,
so $M$ is injective on either side.  Thus $\Sigma_RM$ remains projective
on both sides and represents $M[1]$.

In particular, fix the actual finite bimodule
\begin{equation}
 \label{branch:Y}
 \mathcal Y=\Sigma_R^4C_1.
\end{equation}
It is projective on both sides and represents $\mathcal C[3]$.  This fixes
one target for both lifts of Proposition~\ref{lift:main}.  Concretely, the
tail map there at index $4$ induces
$\Omega_R^4U_H\to C_1$, because
\[
 \coker\bigl(K[3]_5\longrightarrow K[3]_4\bigr)
 =\coker(K_2\longrightarrow K_1).
\]
Taking four cosyzygies gives the stable map $U_H\to \mathcal Y$, using
$\Sigma_R^4\Omega_R^4U_H\cong U_H$ stably.  It therefore has an actual
bimodule-map representative between these finite modules.

For $i=1,2$, put
\[
 \sigma_i=h_{H_i}\otimes h_{H_i},\qquad U_i=A_{\sigma_i}.
\]
Thus $U_i$ has its usual left action and right action $u\cdot a=u\sigma_i(a)$.
Each $U_i$ is free of rank one on both $A$-sides.  The canonical
identification $U_i\otimes_A X=X$ uses the fact that $\sigma_i$ fixes the
simple character of $X$.

By Lemma~\ref{lift:evaluation}, the top projections of the evaluated lifts
are $H_i^2\beta_0$ for one fixed nonzero class
$\beta_0\in\stHom{A}{X}{X[-1]}$.  By
Proposition~\ref{cone:profile}, the top projection
\[
 W^3=\stHom{A}{X}{\mathcal C X[3]}
 \longrightarrow \stHom{A}{X}{X[-1]}
\]
is an isomorphism.  Rescale the two lifts by $H_i^{-2}$, and choose actual
representatives $g_i:U_i\to \mathcal Y$ of these rescaled maps.  Their evaluations
are the same nonzero element $w\in W^3$.

Take the finite free $R$-module $Q=R\otimes_{\kk}\mathcal Y$ and its surjection
$\pi:Q\to \mathcal Y$, $\pi(r\otimes y)=ry$.  Define $F$ by the exact sequence
\begin{equation}
 \label{branch:finite-fiber}
 0\longrightarrow F\longrightarrow U_1\oplus U_2\oplus Q
 \xrightarrow{\ (g_1,g_2,\pi)\ }\mathcal Y\longrightarrow0,
\end{equation}
and let $\rho_i:F\to U_i$ be the two projections.  This is a construction
of finite modules: $C_1$ is a cokernel of two finite terms, four finite
cokernels give $\mathcal Y$, and one finite kernel gives $F$.  Since $\mathcal Y$ is
projective on each side, \eqref{branch:finite-fiber} splits on each side.
Consequently $F$ is projective as both a left and a right $A$-module.

Tensoring \eqref{branch:finite-fiber} with $X$ preserves exactness.
Furthermore, $Q\otimes_AX$ is projective over $A$: evaluation of a free
$A^e$-module is a free $A$-module, and evaluation preserves direct
summands.  Thus the evaluated stable triangle is
\begin{equation}
 \label{branch:evaluated-triangle}
 FX\longrightarrow X\oplus X\xrightarrow{\ (w,w)\ }
 \mathcal C X[3]\longrightarrow FX[1].
\end{equation}
Here and below $FM$ means $F\otimes_AM$.  Tensoring with $F$ is exact
because $F$ is right projective; it preserves projectives because $F$ is
left projective.  It therefore acts on stable maps compatibly with shifts.
The actual maps $\rho_i$ give natural transformations from this functor
to tensoring with $U_i$.

\begin{lemma}
 \label{branch:profile}
 Put $V^a=\stHom{A}{X}{FX[a]}$ for each integer $a$.  There is a unique
 $v\in V^0$ whose two projections are $\id_X$.  It spans $V^0$.
 For every integer $a>0$, projection induces an isomorphism
 \begin{equation}
  \label{branch:positive-projection}
  (\rho_1[a],\rho_2[a]):V^a\xrightarrow{\sim}
  (H^a)^2,
  \qquad H^a=\stHom{A}{X}{X[a]}.
 \end{equation}
\end{lemma}

\begin{proof}
Applying stable Hom from $X$ to \eqref{branch:evaluated-triangle} gives
\[
 (H^{a-1})^2\longrightarrow W^{a+2}
 \longrightarrow V^a\longrightarrow(H^a)^2
 \longrightarrow W^{a+3}.
\]
At $a=0$, the first relevant group is $W^2=0$, and the last map is
$\kk^2\to W^3$, $(c,d)\mapsto(c+d)w$.  It is surjective, so $V^0$
identifies with the diagonal line in $\kk^2$.  This proves the assertion
about $v$.

At $a=1$, the preceding map
$(H^0)^2\to W^3$ is precisely this surjection.  It kills the
potential connecting term from $W^3$ to $V^1$, and $W^4=0$.  Therefore
\eqref{branch:positive-projection} is an isomorphism in degree $1$ as
well.  For $a\ge2$, both $W^{a+2}$ and $W^{a+3}$ vanish by
Proposition~\ref{cone:profile}, giving the isomorphism directly.
\end{proof}

We have obtained an actual finite bimodule $F$ and a nonzero stable map
$v:X\to FX$.  It remains to compare tensoring an extension with $F$ and
composing that extension with $v$.

\subsection{The twist eigenvalues and the comparison map}

\begin{lemma}
 \label{branch:twist}
 For $H\in\{H_1,H_2\}$, tensoring with
 $U_H=A_{h_H\otimes h_H}$, under its canonical identification on $X$,
 acts on the whole space $H^{3m}$ as multiplication by $H^{-m}$
 for every integer $m\ge0$.
\end{lemma}

\begin{proof}
The inverse automorphism in this formula follows from the right-twist
convention.  For an automorphism $h$ of $T$ and a left $T$-module $M$,
the balanced map
\[
 T_h\otimes_TM\longrightarrow {}_{h^{-1}}M,
 \qquad a\otimes m\longmapsto h^{-1}(a)m
\]
is an isomorphism; the target action is $b\cdot m=h^{-1}(b)m$.
For $M=s$ it is the canonical simple identification.

In the normalized bar resolution of $s$, a comparison into its
$T_h$-tensor image applies $h^{-1}$ to every bar letter.  It commutes
with the differential because $h^{-1}$ preserves products; the first
multiplication agrees by the balancing relation for the right-twisted
factor, and the last action agrees because $h$ fixes the character of
$s$.  The simple-valued degree-three cocycle obtained from $p$ has
exactly one dual-ideal input in each nonzero component.  Indeed $p$
lowers the dual-letter weight by one, and only weight-zero outputs can
act nontrivially on $s$.  The bar comparison therefore multiplies this
cochain, and hence its class $\tau$, by $H^{-1}$.

The tensor functor preserves Yoneda products.  In the two-factor algebra
$A$, both factors of one branch use the same $H$, so both $\tau_1$ and
$\tau_2$ acquire the factor $H^{-1}$.  Every basis monomial
$\tau_1^r\tau_2^{m-r}$ of $H^{3m}$ acquires $H^{-m}$.
Linearity gives the assertion on the entire homogeneous piece.
\end{proof}

\begin{proposition}
 \label{branch:comparison}
 The finite bimodule $F$ in \eqref{branch:finite-fiber} is projective on
 both $A$-sides, and the map $v$ of Lemma~\ref{branch:profile} satisfies
 the following conditions.  For every integer $a\ge0$, define
 \begin{equation}
  \label{branch:delta}
  \delta^a:H^a\oplus H^a\longrightarrow V^a,
  \qquad (g,j)\longmapsto F(g)v+v[a]j.
 \end{equation}
 Then $\delta^0$ is surjective with a one-dimensional kernel, and
 $\delta^a$ is an isomorphism for every integer $a>0$.
\end{proposition}

\begin{proof}
Finiteness and both projectivity assertions were proved in the
construction.  The actual bimodule projections $\rho_i$ are natural
with respect to tensoring module maps and their stable shifts.  Since
$\rho_i v=\id_X$, they give
\[
 \rho_i[a]F(g)v=U_i(g),\qquad
 \rho_i[a]v[a]j=j.
\]
For $a=0$, $\kk$-linearity therefore gives
$\delta^0(c\id_X,d\id_X)=(c+d)v$.  This is onto and has the diagonal
line as its kernel.

Let $a=3m>0$.  Under \eqref{branch:positive-projection} and
Lemma~\ref{branch:twist}, the map \eqref{branch:delta} is
\begin{equation}
 \label{branch:matrix}
 (g,j)\longmapsto
 \bigl(H_1^{-m}g+j,\ H_2^{-m}g+j\bigr).
\end{equation}
It is a $2\times2$ scalar matrix tensored with the identity of
$H^{3m}$; that space has dimension $m+1$.  The scalar matrix has determinant
$H_1^{-m}+H_2^{-m}\ne0$.  Otherwise $H_1^m=H_2^m$, contrary to
algebraic independence: the two distinct monomials remain distinct in
the fraction field of $\mathbb F_2[q,H_1,H_2]$.  This argument holds for
every positive integer $m$, including those divisible by the
characteristic.  Hence \eqref{branch:matrix} is invertible.

If $a>0$ is not divisible by $3$, then $H^a=0$, and
Lemma~\ref{branch:profile} gives $V^a=0$.  Thus $\delta^a$ is an
isomorphism in these degrees too.
\end{proof}

The hypotheses of Proposition~\ref{conv:proposition} are now verified for the
finite symmetric algebra $A$, its finite module $X$, and the actual
finite bimodule $F$ and stable map $v$.  Applying that proposition to
\[
 \Lambda=\begin{pmatrix}A&0\\ F&A\end{pmatrix}
\]
produces the required finite nonprojective Gorenstein-projective left
module, with both positive Ext vanishings in every degree.  In
particular, the degree-one assertion uses the surjectivity at degree
zero established in Proposition~\ref{branch:comparison}, in addition to the
positive-degree invertibility.
